\scanpage{423}
\mspage{162}
prouver que cette catégorie $\mathcal C(c,u,v)$ est asphérique. Pour cela, on regarde le chemin composé
\[
\widetilde c=\mathrm{ch}(u)\,c\,\mathrm{ch}(v)^{\unsure{o|0}}
\]
où $\mathrm{ch}(u)$ est le chemin \inserted{(direct)} de $a'$ \unsure{vers} $a$ défini par $u:a'\mathrel{\unsure{\text{---}|\to}}a$, et $\mathrm{ch}(\unsure{u|v})^o$ désigne donc le chemin (rétrograde) de $b$ à $b'$ défini par $v:b'\to b$. \deleted{\unsure{du}} $c$ est de type $\tau$, $\widetilde c$ est donc de type $\Delta_+\tau\Delta_-$.

On va comparer $\mathcal C(c;u,v)$ et la catégorie $\underline{\mathrm{Ch}}(X)_{/\widetilde c}$, à l'aide de deux foncteurs
\diagram{\begin{tikzcd}[ampersand replacement=\&,column sep=large]
\mathcal C(c,u,v)\arrow[r,shift left=.6ex,"\varphi"] \& \underline{\mathrm{Ch}}(X\unsure{,|;}a',b')_{/\widetilde c}\arrow[l,shift left=.6ex,"\psi"],
\end{tikzcd}}
dont on prouvera qu'ils définissent des flèches inverses l'une de l'autre dans Hot. Comme $\underline{\mathrm{Ch}}(X)_{/\widetilde c}$ est contractile (ayant un objet final), donc \deleted{\unsure{Hot}}asphérique, il en est de même de $\mathcal C(c,u,v)$, et on gagne.

\diagram{\begin{tikzcd}[ampersand replacement=\&,row sep=small,column sep=1em]
\&\&\makebox[0pt]{$\overbrace{\hspace{6cm}}^{\widetilde{c'}}$}\&\&\\
a'\arrow[r,"\mathrm{id}"]\arrow[d,"\mathrm{id}"'] \& a'\arrow[rr,dashed,dash,"c'"]\arrow[d,"u"'] \& {}\arrow[d,dashed,"w'"] \& b'\arrow[d,"v"'] \& b'\arrow[l,"\mathrm{id}"']\arrow[d,"\mathrm{id}"]\\
a'\arrow[r,"u"] \& a\arrow[rr,dash,"c"'] \& {} \& b \& b'\arrow[l,"v"']\\
\&\&\makebox[0pt]{$\underbrace{\hspace{6cm}}_{\widetilde c}$}\&\&
\arrow[from=2-1,to=3-2,phantom,"*" description]
\arrow[from=2-4,to=3-5,phantom,"*" description]
\end{tikzcd}}
On définit $\varphi$ par
\[
\boxed{\varphi(c',w')=(\widetilde{c'},\widetilde{w'})}
\]
où ($c'$ étant de type $\tau'$) $\widetilde{c'}$ est le chemin de type $\Delta^+\tau'\Delta^-$ déduit de \deleted{$\unsure{\varphi}$} $c'$ \inserted{\unsure{comme image inv.|comme image inverse|canoniquement}} par la dégénérescence $\Delta^+\tau'\Delta^-\to\tau'$, i.e. on prolonge $c'$ d'un ``cran \inserted{$\Delta_+^{\deleted{\unsure{+}}}$}'' $a'\xrightarrow{\mathrm{id}}a'$ à gauche, et d'un ``cran $\Delta_-^{\deleted{\unsure{-}}}$'' $b'\xleftarrow{\mathrm{id}}b'$ à droite, cf figure.
% archive sheet number: 422
